\documentclass[11pt,a4paper]{article}

\usepackage[utf8]{inputenc}
\usepackage[T1]{fontenc}
\usepackage[french]{babel}
\usepackage[a4paper,margin=2.2cm]{geometry}
\usepackage{amsmath,amssymb,amsthm,mathtools}
\usepackage{lmodern}
\usepackage{enumitem}
\usepackage{hyperref}
\usepackage{xcolor}
\usepackage{fancyhdr}
\usepackage{stmaryrd}
\usepackage{mathrsfs}
\usepackage{array}
\usepackage{bm}

\hypersetup{
    colorlinks=true,
    linkcolor=black,
    urlcolor=blue,
    citecolor=black
}

\setlength{\parindent}{0pt}
\setlength{\parskip}{0.6em}

\pagestyle{fancy}
\fancyhf{}
\fancyhead[L]{Réduction --- Chapitre 16}
\fancyhead[R]{Diagonalisation}
\fancyfoot[C]{\thepage}

% Environnements
\newtheorem{definition}{Définition}[section]
\newtheorem{theoreme}[definition]{Théorème}
\newtheorem{proposition}[definition]{Proposition}
\newtheorem{corollaire}[definition]{Corollaire}
\newtheorem{lemme}[definition]{Lemme}
\theoremstyle{remark}
\newtheorem{remarque}[definition]{Remarque}
\newtheorem{exemple}[definition]{Exemple}
\newtheorem*{methode}{Méthode}
\newtheorem*{idee}{Idée}

% Commandes
\newcommand{\R}{\mathbb{R}}
\newcommand{\C}{\mathbb{C}}
\newcommand{\Q}{\mathbb{Q}}
\newcommand{\Z}{\mathbb{Z}}
\newcommand{\N}{\mathbb{N}}
\newcommand{\K}{\mathbb{K}}
\newcommand{\Vect}{\mathrm{Vect}}
\newcommand{\Ker}{\mathrm{Ker}}
\newcommand{\Img}{\mathrm{Im}}
\newcommand{\rg}{\mathrm{rg}}
\newcommand{\id}{\mathrm{Id}}
\newcommand{\Sp}{\mathrm{Sp}}
\newcommand{\tr}{\mathrm{tr}}
\newcommand{\diag}{\mathrm{diag}}
\newcommand{\Mnn}{M_n(\K)}
\newcommand{\GLn}{GL_n(\K)}
\newcommand{\chiA}{\chi_A}
\newcommand{\chil}{\chi_u}

\begin{document}

\begin{center}
    {\LARGE \textbf{Chapitre 16 --- Diagonalisation}}\\[0.4em]
    {\large Réduire un endomorphisme à sa forme la plus simple}
\end{center}

\hrule
\vspace{1em}

\tableofcontents

\newpage

\section{Pourquoi diagonaliser ?}

Dans l’étude d’un endomorphisme \(u\) ou d’une matrice carrée \(A\), la situation idéale est celle où l’on peut trouver une base dans laquelle la matrice devient diagonale.

Pourquoi ? Parce qu’une matrice diagonale est extrêmement simple à manipuler :
\[
D=
\diag(\lambda_1,\dots,\lambda_n).
\]
Alors
\[
D^p=\diag(\lambda_1^p,\dots,\lambda_n^p),
\]
et plus généralement
\[
P(D)=\diag(P(\lambda_1),\dots,P(\lambda_n))
\]
pour tout polynôme \(P\).

\begin{idee}
Diagonaliser, c’est remplacer un problème global compliqué par des calculs séparés sur chaque direction propre.
\end{idee}

La diagonalisation est donc un outil central pour :
\begin{itemize}
    \item calculer des puissances de matrices ;
    \item étudier des suites récurrentes linéaires ;
    \item comprendre le comportement asymptotique de \(A^n\) ;
    \item préparer le théorème spectral et les réductions plus fines.
\end{itemize}

\section{Définition}

\subsection{Endomorphisme diagonalisable}

\begin{definition}
Soit \(E\) un \(\K\)-espace vectoriel de dimension finie, et
\[
u\in \mathcal L(E,E).
\]
On dit que \(u\) est \textbf{diagonalisable} s’il existe une base de \(E\) formée de vecteurs propres de \(u\).
\end{definition}

\subsection{Matrice diagonalisable}

\begin{definition}
Une matrice \(A\in M_n(\K)\) est dite \textbf{diagonalisable} s’il existe une matrice inversible \(P\in GL_n(\K)\) et une matrice diagonale \(D\in M_n(\K)\) telles que
\[
A=PDP^{-1}.
\]
\end{definition}

\begin{remarque}
Dire que \(A\) est diagonalisable revient à dire qu’elle est semblable à une matrice diagonale.
\end{remarque}

\subsection{Équivalence des deux points de vue}

\begin{proposition}
Un endomorphisme \(u\) est diagonalisable si et seulement si sa matrice dans une base est diagonalisable.
\end{proposition}

\begin{proof}
Si \(u\) est diagonalisable, il existe une base \(\mathcal B\) de vecteurs propres. Dans cette base, la matrice de \(u\) est diagonale.

Réciproquement, si la matrice de \(u\) dans une base \(\mathcal B\) est diagonale, alors les vecteurs de \(\mathcal B\) sont des vecteurs propres de \(u\). Donc \(u\) est diagonalisable.
\end{proof}

\section{Première caractérisation}

\begin{theoreme}
Soit \(u\in \mathcal L(E,E)\). Alors \(u\) est diagonalisable si et seulement si
\[
E=\bigoplus_{\lambda\in\Sp(u)} E_\lambda(u).
\]
\end{theoreme}

\begin{proof}
Supposons \(u\) diagonalisable. Il existe une base \(\mathcal B\) de \(E\) formée de vecteurs propres. En regroupant les vecteurs de \(\mathcal B\) associés à une même valeur propre, on obtient que \(E\) est somme des sous-espaces propres. Cette somme est directe, car des sous-espaces propres associés à des valeurs propres distinctes sont en somme directe.

Réciproquement, si
\[
E=\bigoplus_{\lambda\in\Sp(u)}E_\lambda(u),
\]
on choisit une base de chaque sous-espace propre, puis on réunit ces bases. On obtient une base de \(E\) composée de vecteurs propres. Donc \(u\) est diagonalisable.
\end{proof}

\begin{corollaire}
Un endomorphisme \(u\) est diagonalisable si et seulement si la somme des dimensions de ses sous-espaces propres vaut \(\dim(E)\).
\end{corollaire}

\begin{proof}
Les sous-espaces propres associés à des valeurs propres distinctes sont en somme directe. Donc la dimension de leur somme est la somme de leurs dimensions. Cette somme vaut \(E\) si et seulement si cette somme de dimensions vaut \(\dim(E)\).
\end{proof}

\section{Critère des valeurs propres distinctes}

\begin{theoreme}
Soit \(E\) de dimension \(n\). Si un endomorphisme \(u\) possède \(n\) valeurs propres distinctes dans \(\K\), alors \(u\) est diagonalisable.
\end{theoreme}

\begin{proof}
Pour chaque valeur propre \(\lambda_i\), choisissons un vecteur propre non nul \(x_i\in E_{\lambda_i}\). Les valeurs propres étant deux à deux distinctes, la famille \((x_1,\dots,x_n)\) est libre. Comme \(E\) est de dimension \(n\), c’est une base. Cette base est composée de vecteurs propres, donc \(u\) est diagonalisable.
\end{proof}

\begin{remarque}
C’est le critère le plus rapide à utiliser en pratique.
\end{remarque}

\section{Critère complet par multiplicités}

\subsection{Rappel}

Pour une valeur propre \(\lambda\) :
\begin{itemize}
    \item sa \textbf{multiplicité algébrique} est sa multiplicité comme racine du polynôme caractéristique ;
    \item sa \textbf{multiplicité géométrique} est
    \[
    \dim(E_\lambda).
    \]
\end{itemize}

On sait déjà que
\[
1\le \dim(E_\lambda)\le m_\lambda.
\]

\subsection{Critère}

\begin{theoreme}
Soit \(u\in\mathcal L(E,E)\), avec \(\dim(E)=n\). Alors \(u\) est diagonalisable si et seulement si :
\begin{enumerate}[label=\arabic*)]
    \item le polynôme caractéristique \(\chil\) est scindé sur \(\K\) ;
    \item pour toute valeur propre \(\lambda\),
    \[
    \dim(E_\lambda)=m_\lambda.
    \]
\end{enumerate}
\end{theoreme}

\begin{proof}
Supposons \(u\) diagonalisable. Dans une base de vecteurs propres, la matrice de \(u\) est diagonale :
\[
D=\diag(\lambda_1,\dots,\lambda_n).
\]
Le polynôme caractéristique est donc
\[
\chil(X)=\prod_{i=1}^n (X-\lambda_i),
\]
donc il est scindé. Pour chaque valeur propre \(\lambda\), sa multiplicité algébrique est le nombre de fois où \(\lambda\) apparaît sur la diagonale, et cette quantité est aussi la dimension du sous-espace propre correspondant.

Réciproquement, supposons
\[
\chil(X)=\prod_{i=1}^r (X-\lambda_i)^{m_i}
\]
avec les \(\lambda_i\) distinctes, et
\[
\dim(E_{\lambda_i})=m_i
\quad \text{pour tout } i.
\]
Les sous-espaces propres sont en somme directe, donc
\[
\dim\left(E_{\lambda_1}\oplus\cdots\oplus E_{\lambda_r}\right)
=
\sum_{i=1}^r \dim(E_{\lambda_i})
=
\sum_{i=1}^r m_i
=
n.
\]
Cette somme directe est donc de dimension \(n\), donc égale à \(E\). Ainsi
\[
E=\bigoplus_{i=1}^r E_{\lambda_i},
\]
et \(u\) est diagonalisable.
\end{proof}

\begin{corollaire}
Si une valeur propre \(\lambda\) vérifie
\[
\dim(E_\lambda)<m_\lambda,
\]
alors l’endomorphisme n’est pas diagonalisable.
\end{corollaire}

\section{Cas des matrices}

\begin{theoreme}
Une matrice \(A\in M_n(\K)\) est diagonalisable si et seulement s’il existe une base de \(\K^n\) formée de vecteurs propres de \(A\).
\end{theoreme}

\begin{proof}
C’est exactement la traduction matricielle de la définition d’un endomorphisme diagonalisable.
\end{proof}

\begin{proposition}
Soient \(A\in M_n(\K)\), \(\lambda_1,\dots,\lambda_n\in\K\), et \(v_1,\dots,v_n\in\K^n\) des vecteurs propres associés. Si \((v_1,\dots,v_n)\) est une base, alors la matrice
\[
P=\begin{pmatrix}
| & & |\\
v_1 & \cdots & v_n\\
| & & |
\end{pmatrix}
\]
vérifie
\[
A=PDP^{-1},
\]
où
\[
D=\diag(\lambda_1,\dots,\lambda_n).
\]
\end{proposition}

\begin{proof}
Par définition des vecteurs propres,
\[
Av_j=\lambda_j v_j
\qquad \text{pour } j=1,\dots,n.
\]
Cela s’écrit matriciellement
\[
AP=PD.
\]
Comme \(P\) est inversible (ses colonnes forment une base), on obtient
\[
A=PDP^{-1}.
\]
\end{proof}

\begin{remarque}
Dans une diagonalisation, la matrice \(P\) est constituée de vecteurs propres en colonnes.
\end{remarque}

\section{Comment diagonaliser effectivement une matrice ?}

\begin{methode}
Pour diagonaliser une matrice \(A\in M_n(\K)\) :
\begin{enumerate}[label=\arabic*)]
    \item calculer le polynôme caractéristique \(\chi_A\) ;
    \item déterminer les valeurs propres ;
    \item pour chaque valeur propre \(\lambda\), calculer le sous-espace propre
    \[
    E_\lambda=\Ker(A-\lambda I_n) ;
    \]
    \item vérifier que la somme des dimensions des sous-espaces propres vaut \(n\) ;
    \item choisir une base de chaque sous-espace propre ;
    \item former \(P\) avec les vecteurs propres en colonnes ;
    \item former \(D\) avec les valeurs propres correspondantes sur la diagonale.
\end{enumerate}
\end{methode}

\section{Exemples détaillés}

\subsection{Exemple avec deux valeurs propres distinctes}

\begin{exemple}
Considérons
\[
A=
\begin{pmatrix}
1 & 1\\
0 & 2
\end{pmatrix}.
\]
On calcule
\[
\chi_A(X)=\det\begin{pmatrix}
X-1 & -1\\
0 & X-2
\end{pmatrix}
=(X-1)(X-2).
\]
Les valeurs propres sont \(1\) et \(2\), distinctes, donc \(A\) est diagonalisable.

Cherchons les vecteurs propres.

Pour \(\lambda=1\),
\[
A-I_2=
\begin{pmatrix}
0 & 1\\
0 & 0
\end{pmatrix},
\]
donc
\[
E_1=\left\{
\begin{pmatrix}
x\\0
\end{pmatrix}
: x\in\R
\right\}
=
\Vect\left(
\begin{pmatrix}
1\\0
\end{pmatrix}
\right).
\]

Pour \(\lambda=2\),
\[
A-2I_2=
\begin{pmatrix}
-1 & 1\\
0 & 0
\end{pmatrix},
\]
donc
\[
E_2=\left\{
\begin{pmatrix}
x\\x
\end{pmatrix}
: x\in\R
\right\}
=
\Vect\left(
\begin{pmatrix}
1\\1
\end{pmatrix}
\right).
\]

On prend
\[
P=
\begin{pmatrix}
1 & 1\\
0 & 1
\end{pmatrix},
\qquad
D=
\begin{pmatrix}
1 & 0\\
0 & 2
\end{pmatrix}.
\]
Alors
\[
A=PDP^{-1}.
\]
\end{exemple}

\subsection{Exemple non diagonalisable}

\begin{exemple}
Considérons
\[
A=
\begin{pmatrix}
1 & 1\\
0 & 1
\end{pmatrix}.
\]
On a
\[
\chi_A(X)=(X-1)^2.
\]
L’unique valeur propre est \(1\), de multiplicité algébrique \(2\).

Or
\[
A-I_2=
\begin{pmatrix}
0 & 1\\
0 & 0
\end{pmatrix},
\]
d’où
\[
E_1=
\left\{
\begin{pmatrix}
x\\0
\end{pmatrix}
: x\in\R
\right\}.
\]
Donc
\[
\dim(E_1)=1<2.
\]
La matrice n’est pas diagonalisable.
\end{exemple}

\subsection{Exemple avec multiplicité algébrique \(>1\) et diagonalisabilité}

\begin{exemple}
Considérons
\[
A=
\begin{pmatrix}
2 & 0 & 0\\
0 & 3 & 0\\
0 & 0 & 3
\end{pmatrix}.
\]
Le polynôme caractéristique est
\[
\chi_A(X)=(X-2)(X-3)^2.
\]
La valeur propre \(3\) a multiplicité algébrique \(2\), et son sous-espace propre est
\[
E_3=
\Vect\left(
\begin{pmatrix}
0\\1\\0
\end{pmatrix},
\begin{pmatrix}
0\\0\\1
\end{pmatrix}
\right),
\]
de dimension \(2\). Donc \(A\) est diagonalisable.
\end{exemple}

\section{Diagonalisation et calculs de puissances}

\subsection{Puissances}

\begin{theoreme}
Si
\[
A=PDP^{-1}
\]
avec \(D\) diagonale, alors pour tout \(n\in\N\),
\[
A^n=PD^nP^{-1}.
\]
\end{theoreme}

\begin{proof}
Par récurrence ou en utilisant l’associativité du produit matriciel :
\[
A^2=(PDP^{-1})(PDP^{-1})=PD^2P^{-1},
\]
et de même pour tout \(n\).
\end{proof}

\begin{proposition}
Si
\[
D=\diag(\lambda_1,\dots,\lambda_n),
\]
alors
\[
D^p=\diag(\lambda_1^p,\dots,\lambda_n^p).
\]
\end{proposition}

\begin{proof}
Le calcul se fait coefficient par coefficient sur la diagonale.
\end{proof}

\begin{remarque}
C’est l’une des principales utilités de la diagonalisation.
\end{remarque}

\subsection{Exemple}

\begin{exemple}
Avec
\[
A=
\begin{pmatrix}
1 & 1\\
0 & 2
\end{pmatrix}
=
PDP^{-1},
\qquad
D=
\begin{pmatrix}
1 & 0\\
0 & 2
\end{pmatrix},
\]
on obtient
\[
A^n=PD^nP^{-1}
=
P
\begin{pmatrix}
1 & 0\\
0 & 2^n
\end{pmatrix}
P^{-1}.
\]
\end{exemple}

\section{Diagonalisation et polynômes de matrices}

\begin{theoreme}
Si
\[
A=PDP^{-1}
\]
et \(Q\in\K[X]\), alors
\[
Q(A)=PQ(D)P^{-1}.
\]
\end{theoreme}

\begin{proof}
Si
\[
Q(X)=a_0+a_1X+\cdots+a_mX^m,
\]
alors
\[
Q(A)=a_0I_n+a_1A+\cdots+a_mA^m.
\]
Or
\[
A^k=PD^kP^{-1},
\]
donc
\[
Q(A)=P(a_0I_n+a_1D+\cdots+a_mD^m)P^{-1}=PQ(D)P^{-1}.
\]
\end{proof}

\begin{corollaire}
Si
\[
D=\diag(\lambda_1,\dots,\lambda_n),
\]
alors
\[
Q(D)=\diag(Q(\lambda_1),\dots,Q(\lambda_n)).
\]
\end{corollaire}

\begin{proof}
On applique le polynôme \(Q\) terme à terme sur la diagonale.
\end{proof}

\section{Cas particuliers importants}

\subsection{Projecteurs}

\begin{proposition}
Tout projecteur \(p\) vérifiant
\[
p^2=p
\]
est diagonalisable.
\end{proposition}

\begin{proof}
Le polynôme
\[
X^2-X=X(X-1)
\]
annule \(p\). Il est scindé à racines simples. Donc \(p\) est diagonalisable.

Plus explicitement, les seules valeurs propres possibles sont \(0\) et \(1\), et on a
\[
E=\Ker(p)\oplus \Img(p).
\]
\end{proof}

\subsection{Involutions}

\begin{proposition}
Si \(u^2=\id_E\) et si \(\mathrm{car}(\K)\neq 2\), alors \(u\) est diagonalisable.
\end{proposition}

\begin{proof}
Le polynôme
\[
X^2-1=(X-1)(X+1)
\]
annule \(u\). Il est scindé à racines simples, donc \(u\) est diagonalisable.
\end{proof}

\subsection{Matrices symétriques réelles}

\begin{remarque}
Les matrices symétriques réelles sont diagonalisables, et même orthogonalement diagonalisables. Cela sera étudié précisément dans le chapitre sur la réduction des endomorphismes symétriques.
\end{remarque}

\section{Lien avec les suites récurrentes linéaires}

\subsection{Principe}

Les suites récurrentes linéaires se réécrivent souvent sous forme matricielle :
\[
X_{n+1}=AX_n.
\]
Alors
\[
X_n=A^nX_0.
\]
Si \(A\) est diagonalisable, le calcul de \(A^n\) devient explicite.

\subsection{Exemple de Fibonacci}

\begin{exemple}
La suite de Fibonacci vérifie
\[
F_{n+2}=F_{n+1}+F_n.
\]
Posons
\[
X_n=
\begin{pmatrix}
F_{n+1}\\
F_n
\end{pmatrix}.
\]
Alors
\[
X_{n+1}=
\begin{pmatrix}
1 & 1\\
1 & 0
\end{pmatrix}
X_n.
\]
Notons
\[
A=
\begin{pmatrix}
1 & 1\\
1 & 0
\end{pmatrix}.
\]
Son polynôme caractéristique est
\[
\chi_A(X)=X^2-X-1.
\]
Ses valeurs propres sont
\[
\varphi=\frac{1+\sqrt5}{2},
\qquad
\psi=\frac{1-\sqrt5}{2},
\]
distinctes, donc \(A\) est diagonalisable. On peut alors exprimer \(A^n\), puis obtenir une formule explicite pour \(F_n\).
\end{exemple}

\begin{remarque}
C’est l’une des applications classiques et importantes de la diagonalisation.
\end{remarque}

\section{Diagonalisation et décomposition spectrale}

\begin{proposition}
Si \(u\) est diagonalisable et admet pour valeurs propres distinctes \(\lambda_1,\dots,\lambda_r\), alors
\[
E=\bigoplus_{i=1}^r E_{\lambda_i}.
\]
\end{proposition}

\begin{remarque}
La diagonalisation consiste alors à choisir une base adaptée à cette décomposition.
\end{remarque}

\section{Critères utiles à retenir}

\begin{theoreme}
Pour \(u\in\mathcal L(E,E)\), les assertions suivantes sont équivalentes :
\begin{enumerate}[label=\arabic*)]
    \item \(u\) est diagonalisable ;
    \item il existe une base de \(E\) formée de vecteurs propres ;
    \item
    \[
    E=\bigoplus_{\lambda\in\Sp(u)}E_\lambda ;
    \]
    \item le polynôme caractéristique est scindé et, pour toute valeur propre \(\lambda\),
    \[
    \dim(E_\lambda)=m_\lambda.
    \]
\end{enumerate}
\end{theoreme}

\begin{proof}
Les équivalences ont été démontrées dans les sections précédentes.
\end{proof}

\section{Méthode pratique complète}

\begin{methode}
Pour diagonaliser un endomorphisme ou une matrice :
\begin{enumerate}[label=\arabic*)]
    \item calculer le polynôme caractéristique ;
    \item déterminer les valeurs propres ;
    \item calculer les sous-espaces propres ;
    \item vérifier la diagonalisabilité ;
    \item construire une base de vecteurs propres ;
    \item former \(P\) avec ces vecteurs en colonnes ;
    \item former \(D\) avec les valeurs propres associées ;
    \item écrire
    \[
    A=PDP^{-1}.
    \]
\end{enumerate}
\end{methode}

\section{Erreurs classiques}

\begin{itemize}
    \item Croire qu’un polynôme caractéristique scindé suffit à assurer la diagonalisabilité.
    \item Confondre multiplicité algébrique et multiplicité géométrique.
    \item Oublier que les colonnes de \(P\) doivent être des vecteurs propres dans le même ordre que les valeurs propres placées dans \(D\).
    \item Se tromper dans l’écriture
    \[
    A=PDP^{-1}
    \]
    en inversant \(P\) et \(P^{-1}\).
    \item Oublier qu’une matrice peut être diagonalisable même avec une valeur propre multiple.
    \item Oublier qu’une matrice triangulaire n’est pas forcément diagonalisable.
\end{itemize}

\section{Résumé du chapitre}

\begin{itemize}
    \item Un endomorphisme est diagonalisable s’il existe une base de vecteurs propres.
    \item Une matrice est diagonalisable si elle est semblable à une matrice diagonale.
    \item Si \(A=PDP^{-1}\), alors
    \[
    A^n=PD^nP^{-1}.
    \]
    \item Si \(A\) a \(n\) valeurs propres distinctes, alors elle est diagonalisable.
    \item Plus généralement, \(A\) est diagonalisable si et seulement si son polynôme caractéristique est scindé et si, pour chaque valeur propre, multiplicité géométrique et multiplicité algébrique coïncident.
    \item Les projecteurs et les involutions sont diagonalisables.
    \item La diagonalisation est un outil fondamental pour calculer des puissances de matrices et résoudre des suites récurrentes linéaires.
\end{itemize}

\section*{Exercices d’application immédiate}

\begin{enumerate}[label=\textbf{Exercice \arabic*.}]
    \item Diagonaliser
    \[
    \begin{pmatrix}
    2 & 0\\
    0 & 3
    \end{pmatrix}.
    \]

    \item Étudier la diagonalisabilité de
    \[
    \begin{pmatrix}
    1 & 1\\
    0 & 2
    \end{pmatrix}.
    \]

    \item Étudier la diagonalisabilité de
    \[
    \begin{pmatrix}
    1 & 1\\
    0 & 1
    \end{pmatrix}.
    \]

    \item Montrer qu’un projecteur est diagonalisable.

    \item Montrer qu’une involution est diagonalisable si \(\mathrm{car}(\K)\neq 2\).

    \item Calculer \(A^n\) pour
    \[
    A=
    \begin{pmatrix}
    1 & 1\\
    0 & 2
    \end{pmatrix}.
    \]

    \item Soit
    \[
    A=
    \begin{pmatrix}
    2 & 1 & 0\\
    0 & 2 & 0\\
    0 & 0 & 3
    \end{pmatrix}.
    \]
    Étudier sa diagonalisabilité.

    \item Soit
    \[
    A=
    \begin{pmatrix}
    1 & 1\\
    1 & 0
    \end{pmatrix}.
    \]
    Diagonaliser \(A\) puis calculer \(A^n\).

    \item Soit \(u\) un endomorphisme dont le polynôme caractéristique est
    \[
    (X-1)^2(X-3),
    \]
    et tel que
    \[
    \dim(E_1)=2,\qquad \dim(E_3)=1.
    \]
    Montrer que \(u\) est diagonalisable.

    \item Soit \(u\) un endomorphisme annulé par
    \[
    X(X-1)(X+2).
    \]
    Montrer que \(u\) est diagonalisable.
\end{enumerate}

\end{document}